\section{Introduction}

% (1) {Overview of IO memory protection}, IOMMU, providing IO memory protection in VMs, vIOMMU to provide intra-guest protection...

% (2) {Performance degradation with intra-guest protection.} Intra-guest IO memory protection has significant performance overheads. We observe up to 90\% throughput degradation (compared to only having inter-guest protection enabled) when using strict intra-guest memory protection. Slightly relaxed when using deferred memory protection (50\% degradation)---at the cost of weaker safety guarantees. \leshna{we don't see degradation in lazy nested. If we focus only on nested than this is a bit deceiving}
% \schaic{When running nested translation, we only see lazy's degradation when VM is running as transmitter.}

(3) Under strict protection, existing Linux network stack performs invalidation upon every unmap operation synchronously by any network application thread. Using {vIOMMU}, processing each invalidation request (IR) within each guest VM involves two sub-operations: (i) IR-PREP---preparing an invalidation request using corresponding IOVA range, and enqueing the request to a (per-IO-device) submission queue and \leshna{these are not per-IO device but per IOMMU (these are different things)}(ii) IR-SUB---submitting the prepared requests from the submission queue(s) to a (per-guest) emulated IOMMU invalidation queue.
\leshna{it is confusing of how many per device queues we have from this statement, looks like all devices have a single per-guest queue, which is not the case}
% results in (ENQ) enqueue of invalidation request to a per-guest emulated invalidation queue, followed by a (FLUSH) flush queue command from the guest 
IR-SUB triggers a VM exit to the hypervisor, which then copies the request to hardware invalidation queue, which invalidates the corresponding IOVA enties in the request from the IOTLB and/or IOPT-caches~\cite{Rubin2024FastSafe}. 


% - default: unmap, enq and flush are synchronous within each vCPU; enq and flush together are serialized across all vCPUs


(4) In existing stacks, abovementioned IR operations provide two interesting properties: (P1) IR-PREPs for any given IO device are executed in a {serialized manner across all threads in a guest VM}. This is done to ensure correctness while performing concurrent writes from all threads performing unmaps to the shared per-IO-device submission queue\leshna{same comment as above, also the serialization is also across a group of IOMMU devices}; and achieved by executing IR-PREP after acquiring a lock on a per-IO-device basis. (P2) Both IR-PREP and IR-SUB operations are {\it coupled together}, i.e., both are executed within the same critical section of the per-IO-device \schai{interrupt-disabling} lock. \saksham{To discuss: what would be the original intended purpose for this coupled design decision? looks like it would minimize the time between IR-PREP and IR-SUB for any IR, basically the request processing delay.}
\schai{The interrupt is disabled to minize processing delay of IR, and the lock is held because IR-SUB requires copying entries from per-device submission queue to emulated queue.}
% IR-PREP operations are serialized across all threads performing unmaps for DMAs from any given IO device 

(5) Fundamentally, P1 and P2 ocurring together is the core reason for aforementioned throughput degradation with existing stacks. 
% Since IR-PREP and IR-SUB are coupled and serialized across all network threads, the aggregate rate of unmap calls across an entire guest VM is bounded, and independent of the number of concurrent threads. 
Since IR-PREP and IR-SUB are coupled and serialized across all network threads, the aggregate rate of processing IRs across all threads is bounded by $R = $1/(time to execute IR-PREP \& IR-SUB) since we can only submit a single invalidation request to the IOMMU HW at a time for processing. Since IR is performed synchronously with unmaps, aggregate rate of unmaps is also bounded by $R$. $R$ is primarily decided by the time to execute IR-SUB, which is orders of magnitude larger than IR-PREP. The IR-SUB operation is much more costly due to involved (unavoidable) VM exits; VM exits are necessary because the guest doesn't have the write privileges to access the IOMMU HW queue. VM exits latencies are typically few tens of microseconds (L = xy microseconds on our setup; \S\ref{sec:inefficiencies}).

Further, using strict protection mode enforces a lower bound on the number of unmap operations required per unit of network achievable throughput ($U = xy$ unmaps per unit network throughput, see \S\ref{sec:inefficiencies} for details). \leshna{this is hard to claim with page pools.} This leads to an upper bound on the achievable network throughput using existing memory protection mechanisms in Linux (provide formula for throughput T as function of U and L here...), which is significantly smaller than the NIC line rates used in our setup (400Gbps). Since the VM exit costs have remained largely stagnant, the network bandwidth utilization is going to get even worse with larger NIC line rates. 
% on average each unmap call will take at least duration required for IR-PREP and IR-SUB operations (illustrated in Figure 2). 

% , \leshna{as mentioned offline, this formula will need to account for a fraction that is page pool miss (to match the observed values), the theoretical value is very small (similar to ones seen in TX)}

% however, if they were decoupled, we can 


% - delegated inval: unmap, enq and flush are synchronous within each vCPU; only flush is serialized across all VCPUs, but batched
(6) We resolve this problem by redesigning the existing memory protection mechanisms to perform concurrent processing of IRs, without impacting strict safety guarantees. We first argue that serializing IR-PREP isn't fundamentally necessary for correctness---instead of using the existing single per-IO-device submission queue for queueing requests from IR-PREP operation, we propose incorporating per-vCPU virtual submission queues, with enqueue to any of these virtual queues requiring a per-vCPU lock. This would enable threads from multiple vCPUs to concurrently execute IR-PREP operation, potentially allowing larger number of IRs to be submitted concurrently to the hardware for processing. 

However, enabling such a concurrent IR processing requires solving an additional challenge: IR-SUB operations in existing stacks also need to be serialized across all threads in the VM to ensure write-safety guarantees for the shared emulated IOMMU invalidation queue. Simply extending the above approach, of using a per-vCPU virtual queue, for the case of IR-SUB is not ideal either. This is because...\saksham{we need to give a good reason here; one argument is it would introduce additional VM exits, another could be it requires modifying guest-hypervisor interface, so is it paravirtualization, and not desirable?} 
\leshna{Do you mean a per-vCPU and per-VM queues? We would need the hypervisor to then serialize them, until there is hardware support for multiple queues, that would just shift serialization down the stack, with more complexity}\saksham{But we might still be able to hide the VM exit overhead by multiple vCPUs triggering VM exits in parallel, right, and conceptually approach the baremetal performance. The question is then, is "added hypervisor complexity" a good enough reason to not adopt this approach?}
% IR-SUB operations are  serialized across all threads performing unmaps for any IOMMU. This is because vIOMMU, designed for unmodified guests, provides the same interface between the guest OS and vIOMMU's emulated invalidation queue, as exists between the host OS and the IOMMU HW invalidation queue: IOMMUs in modern servers typically only support a limited number of invalidation queues (for example, on Intel servers there is a single invalidation queue). Therefore all enqueues to a single HW queue need to be serialized via a lock (to avoid corruption by multiple concurrent enqueues). Hence, viommu also currently uses a lock to serialize enqueue of invalidation request to the emulated invalidation queue across all vCPUs. 
To address this challenge, we propose decoupling the IR-PREP and IR-SUB operations for IR processing. Such a decoupling can be beneficial for concurrent IR processing due to the fact that existing IOMMU invalidation queue interface also supports an enqueue of a batch of invalidation requests (and not just a single request). If IR-PREP operations are not serialized across the vCPUs, we can batch the multiple enqueued requests from multiple vCPUs and perform IR-SUB operation using this batch of requests whenever possible (since IR-SUB needs to be serialized). Even though multiple vCPUs can concurrently execute IRs, IR-PREP and IR-SUB continue to be executed synchronously with unmaps within each thread, hence, the proposed design continues to provide strict protection guarantees. We provide a formal proof in \S\ref{sec:security}. 
\schaic{For bullet point 6, I think we should reverse the two para: introduce decoupling first then per-cpu queue. The reason is that decoupling is the most critical parts which enables batched processing of requests, per-cpu then becomes a natural optimization on top of that. WIP to develop a draft.}\leshna{I agree with siyuan, this story is much cleaner if we introduce decoupling and thus batching first and then the per-CPU idea becomes much easier to explain and a natural progression}
% \saksham{we should probably rename the first design idea to batched invalidation, since it seems pretty important in terms of realizing the benefits?}
% We show in \S\ref{sec:design} that this proposed idea of concurrent IR processing via batched invalidations 

% As long as they both (A) continue to remain synchronous with unmap call within each thread, they continue to provide exactly the same strict protection guarantees; 
% and (B) only FLUSH operations are serialized across all vCPU cores, they continue to provide correct write semantics for the IOMMU HW invalidation queue. 

(7) The above approach of enables scaling network throughput with increasing number of vCPUs (without impacting correctness guarantees). This is because it enables multiple concurrent unmap operations, from multiple vCPUs within the guest, to proceed within an IR-PREP and IR-SUB duration (see Figure 3). Concretely, this enables $N\times$ concurrent requests per duration (where $N$ is the number of vCPUs), instead of only one request from the entire guest in existing design. Implementing this approach---(referred to as \saksham{how about "Batched concurrent invalidations"?})\leshna{this misses the blocking safety property, also don't batched and concurrent here describe the same operation of batching close by requests?}---provided us throughput \todo{xyz}, which is \todo{abc}$\times$ higher than default vIOMMU; enables scaling with vCPUs. The gap between the no vIOMMU performance and the above scheme for any given number of vCPUs, is because of the fact IR-SUB is executed synchronously with IR-PREP. This limits the rate of unmap operations by each vCPU to a maximum of one per IR-PREP+IR-SUB duration; and given the large IR-SUB latency overhead due to VM exits, keeping one request per vCPU over this duration isn't enough to hide the latency, especially with large network bandwidths. 


(8) We argue this problem is not fundamental: \textit{contrary to conventional wisdom, executing both IR-PREP and IR-SUB synchronously with unmap calls are in fact NOT fundamentally necessary to achieve the strict memory protection.} This tradition fallacy arises due to the way existing network drivers recycle physical physical pages after unmap operations, to be mapped to other IO virtual addresses. The drivers today free a physical page after every unmap call. To ensure strict protection, we need to invalidate the IOTLB before the physical page is freed and recycled (or else it could result in known confidentiality attacks~\cite{}). Therefore, today invalidations are performed synchronously with unmaps to ensure they are always executed before the corresponding page free operations. Our key insight is that \textit{as long as the following ordering of operations---unmap$\rightarrow$invalidate$\rightarrow$free page is followed, we maintain the default strict protection guarantees}, irrespective of the fact whether invalidations are synchronous with the unmaps. \saksham{probably need to mention somewhere that recent linux kernel designs also ensure that there is no possibility of TOCTOU attacks even with asynchronous unmaps and invalidations.} \leshna{does this need to go in the intro? I think we should keep the insight of order and refer answers to all questions about securty details to later section}

% - deferred free page: only unmap and enq are synchronous, and flush is asyncrhonous+batched within each vCPU; flush is seialized across vCPUs, but batched
(9) We solve the problem of limited per-vCPU throughput by leveraging the above insight. We propose two architectural modifications: (i) IR-PREP and IR-SUB are executed asynchronously---the thread only executes IR-PREP synchronously with each unmap operation, and defers IR-SUB (again, using a single batched call for all unmap calls) to whenever earliest possible, given the required serialization across all vCPUs; and (ii) page free operations for the set of unmap calls are also deferred until the corresponding IR-SUB has been performed. This ensures that we maintain the required ordering guarantees to ensure strict memory protection guarantees (formal proof in \S\ref{sec:security}). But additionally, it enables each vCPU core to maintain much larger number of parallel unmap calls per unit time, closing the performance gap compared to no-vIOMMU case.


(10) overview of implementation (\saksham{do we propose a "patched guest design"?; need to mention somewhere, we are not paravirtualized; at least we don't seem to be introducing any hypercalls}) and results; close-to-no-viommu results with strict protection guarantees.

need to mention somewhere why not sidecore based approach, what's the qualitative problem with that; it could have eliminated vm exit costs






